We have compared an ytterbium lattice clock in Wuhan and a strontium
lattice clock in Shanghai over a \fibreKm\,km fiber link and
determined the Yb/Sr frequency ratio with a total standard
uncertainty of \uTotalShort$\times10^{-18}$. The value agrees with the
two existing Yb/Sr measurements at the level demanded by the roadmap
for the redefinition of the SI second, and is, to our knowledge, the
first intercity, cross-species frequency-ratio measurement to meet
that requirement. The comparison shows that the limiting factors are
no longer the transfer infrastructure but the systematic evaluations
of the clocks and the geodetic treatment of the gravitational
potential difference between the sites. In particular, the
time-varying gravitational redshift cannot yet be modeled well enough
to correct: the available tide series degrades, rather than improves,
the segment agreement, so the term is carried conservatively as an
uncertainty item. Together with the published same-campus and
continental-network ratios, this measurement marks the convergence of
independent Yb/Sr determinations at the few-$10^{-18}$ level and
demonstrates a working node pair of a future optical-clock network
that will combine fiber backbones with satellite and free-space links.
